\chapter{随机变量的数字特征}

\intro{
	分布函数完整地描述了随机变量的统计规律，但在实际应用中，往往只需要少数几个数字特征就能把握其主要性质。
}

\section{数学期望（均值）}

\begin{mydefinition}{离散型随机变量的期望}
	设离散型随机变量 $X$ 的分布律为 $P\{X = x_k\} = p_k$，若级数 $\sum_k \abs{x_k} p_k$ 收敛，则定义
	\[E(X) = \sum_{k=1}^{\infty} x_k p_k\]
	为 $X$ 的\textbf{数学期望}。
\end{mydefinition}

\begin{mydefinition}{连续型随机变量的期望}
	设连续型随机变量 $X$ 的概率密度函数为 $f(x)$，若 $\int_{-\infty}^{+\infty} \abs{x} f(x)\,dx$ 收敛，则定义
	\[E(X) = \int_{-\infty}^{+\infty} x f(x)\,dx\]
	为 $X$ 的\textbf{数学期望}。
\end{mydefinition}

\begin{remark}
随机变量期望存在的条件是绝对收敛：$\sum \abs{x_k} p_k < \infty$（离散型）或 $\int \abs{x} f(x)\,dx < \infty$（连续型）。
这是为了确保期望值不依赖于求和或积分的次序。Cauchy 分布就是期望不存在的典型例子。
\end{remark}

\begin{theorem}[随机变量函数的期望（LOTUS 定理）]
\begin{itemize}
	\item 一维离散型：$\displaystyle E[g(X)] = \sum_{k} g(x_k) p_k$；
	\item 一维连续型：$\displaystyle E[g(X)] = \int_{-\infty}^{+\infty} g(x) f(x)\,dx$；
	\item 二维连续型：$\displaystyle E[g(X, Y)] = \iint_{\R^2} g(x, y) f(x, y)\,dx\,dy$。
\end{itemize}
该定理的实用价值在于：求 $g(X)$ 的期望时，\textbf{不需要先求出 $Y = g(X)$ 的分布}，直接利用 $X$ 的分布即可。
\end{theorem}

\begin{myproperty}{期望的基本性质}
\begin{enumerate}
	\item $E(C) = C$（常数的期望等于本身）；
	\item \textbf{线性性：} $E(aX + bY + c) = aE(X) + bE(Y) + c$（\textbf{不要求独立！}）；
	\item 若 $X \indep Y$，则 $E(XY) = E(X) E(Y)$（逆命题不成立）；
	\item \textbf{单调性：} 若 $X \leq Y$ a.s.，则 $E(X) \leq E(Y)$；
	\item $\abs{E(X)} \leq E(\abs{X})$。
\end{enumerate}
\end{myproperty}

\subsection{常见分布的期望}

\begin{center}
\begin{tabular}{lcc}
\toprule
\textbf{分布} & \textbf{符号} & \textbf{期望 $E(X)$} \\
\midrule
两点分布 & $B(1, p)$          & $p$ \\
二项分布 & $B(n, p)$          & $np$ \\
Poisson分布 & $P(\lambda)$   & $\lambda$ \\
均匀分布 & $U(a, b)$          & $\dfrac{a+b}{2}$ \\
指数分布 & $E(\lambda)$       & $\dfrac{1}{\lambda}$ \\
正态分布 & $N(\mu, \sigma^2)$ & $\mu$ \\
\bottomrule
\end{tabular}
\end{center}

\section{方差}

\begin{mydefinition}{方差与标准差}
	设 $X$ 为随机变量，其\textbf{方差}定义为：
	\[D(X) = \Var(X) = E\bigl[(X - E(X))^2\bigr].\]
	\textbf{标准差}：$\sigma(X) = \sqrt{D(X)}$。方差刻画了随机变量取值相对于均值的离散程度。
\end{mydefinition}

\begin{theorem}[方差的计算公式]
\[D(X) = E(X^2) - [E(X)]^2.\]
此公式通常比定义式更便于计算。
\end{theorem}

\begin{myproperty}{方差的性质}
\begin{enumerate}
	\item $D(C) = 0$；若 $D(X) = 0$，则 $P\{X = C\} = 1$（退化的随机变量）；
	\item $D(aX + b) = a^2 D(X)$（方差不满足线性性！）；
	\item 若 $X \indep Y$，则 $D(X + Y) = D(X) + D(Y)$；
	\item 一般情形：$D(X \pm Y) = D(X) + D(Y) \pm 2\Cov(X, Y)$。
\end{enumerate}
\end{myproperty}

\subsection{常见分布的方差}

\begin{center}
\begin{tabular}{lcc}
\toprule
\textbf{分布} & \textbf{符号} & \textbf{方差 $D(X)$} \\
\midrule
两点分布 & $B(1, p)$          & $p(1-p)$ \\
二项分布 & $B(n, p)$          & $np(1-p)$ \\
Poisson分布 & $P(\lambda)$   & $\lambda$ \\
均匀分布 & $U(a, b)$          & $\dfrac{(b-a)^2}{12}$ \\
指数分布 & $E(\lambda)$       & $\dfrac{1}{\lambda^2}$ \\
正态分布 & $N(\mu, \sigma^2)$ & $\sigma^2$ \\
\bottomrule
\end{tabular}
\end{center}

\begin{remark}[Poisson 分布的期望与方差相等]
$\lambda$ 既是 Poisson 分布的期望，也是方差，这是一个重要的特征性质。在统计分析中，若样本均值与样本方差接近，则数据可能服从 Poisson 分布。
\end{remark}

\section{Chebyshev 不等式}

\begin{theorem}[Chebyshev 不等式]
设随机变量 $X$ 具有数学期望 $E(X) = \mu$ 和方差 $D(X) = \sigma^2$。则对任意 $\varepsilon > 0$，有
\[P\{\abs{X - \mu} \geq \varepsilon\} \leq \frac{\sigma^2}{\varepsilon^2}.\]
等价形式（令 $\varepsilon = k\sigma$）：
\[P\{\abs{X - \mu} \geq k\sigma\} \leq \frac{1}{k^2}.\]
\end{theorem}

该不等式的威力在于：\textbf{无论随机变量服从什么分布}（只要方差存在），它都成立。它给出了随机变量偏离期望的概率上界。
\begin{itemize}
	\item 取 $k = 2$：偏离 $2\sigma$ 以上的概率不超过 $1/4 = 25\%$；
	\item 取 $k = 3$：偏离 $3\sigma$ 以上的概率不超过 $1/9 \approx 11.1\%$。
\end{itemize}

\begin{remark}
对于正态分布，$P(\abs{X - \mu} \geq 3\sigma) \approx 0.27\%$，远小于 Chebyshev 不等式给出的 $11\%$。这说明 Chebyshev 不等式是保守的界限，但它的普适性使得它在理论证明中非常有用。
\end{remark}

\section{协方差与相关系数}

\begin{mydefinition}{协方差}
	随机变量 $X$ 与 $Y$ 的\textbf{协方差}定义为：
	\[\Cov(X, Y) = E\bigl[(X - E(X))(Y - E(Y))\bigr].\]
	\textbf{计算公式：}
	\[\Cov(X, Y) = E(XY) - E(X)E(Y).\]
\end{mydefinition}

\begin{mydefinition}{不相关}
	若 $\Cov(X, Y) = 0$，则称 $X$ 与 $Y$ \textbf{不相关}。
\end{mydefinition}

\begin{myproperty}{协方差的性质}
\begin{enumerate}
	\item $\Cov(X, X) = D(X)$；
	\item $\Cov(X, Y) = \Cov(Y, X)$（对称性）；
	\item $\Cov(aX, bY) = ab \cdot \Cov(X, Y)$；
	\item $\Cov(X_1 + X_2, Y) = \Cov(X_1, Y) + \Cov(X_2, Y)$（双线性性）；
	\item 若 $X \indep Y$，则 $\Cov(X, Y) = 0$（逆命题不成立，正态分布除外）。
\end{enumerate}
\end{myproperty}

\begin{mydefinition}{相关系数}
	$X$ 与 $Y$ 的\textbf{相关系数}（Pearson correlation coefficient）定义为：
	\[\rho_{XY} = \frac{\Cov(X, Y)}{\sqrt{D(X) D(Y)}} = \frac{\Cov(X, Y)}{\sigma_X \sigma_Y}.\]
\end{mydefinition}

\begin{myproperty}{相关系数的性质}
\begin{enumerate}
	\item $-1 \leq \rho_{XY} \leq 1$（可由 Cauchy--Schwarz 不等式证明）；
	\item $\abs{\rho_{XY}} = 1 \iff X$ 与 $Y$ 以概率 $1$ 存在线性关系 $Y = aX + b$；
	\item $|\rho_{XY}|$ 越接近 $1$，线性相关程度越高；越接近 $0$，线性相关程度越低；
	\item \textbf{对二维正态分布：} 不相关 $\iff$ 独立 $\iff$ $\rho = 0$。
\end{enumerate}
\end{myproperty}

\section{协方差矩阵}

对于 $n$ 维随机变量 $(X_1, \ldots, X_n)$，协方差矩阵 $\Sigma = (\sigma_{ij})_{n \times n}$，其中
\[\sigma_{ij} = \Cov(X_i, X_j) = E[(X_i - E(X_i))(X_j - E(X_j))].\]
协方差矩阵的性质：
\begin{itemize}
	\item 对称矩阵：$\sigma_{ij} = \sigma_{ji}$；
	\item 半正定矩阵：对任意向量 $\mathbf{a}$，$\mathbf{a}^{\mathsf{T}} \Sigma \mathbf{a} \geq 0$；
	\item 对角元素为方差：$\sigma_{ii} = D(X_i)$。
\end{itemize}

\begin{myTheorem}{$n$ 维正态分布的重要特性}
服从 $n$ 维正态分布的随机变量，其两两不相关性与相互独立性是\textbf{等价的}。这是正态分布特有的优良性质。
\end{myTheorem}

\section{独立与不相关的关系}

\begin{myTheorem}{独立与不相关的关系总结}
\begin{enumerate}
	\item 独立 $\Rightarrow$ 不相关（由 $E(XY) = E(X)E(Y)$ 直接得到 $\Cov = 0$）；
	\item 不相关 $\nRightarrow$ 独立（不相关仅意味着没有\textbf{线性关系}，但仍可能存在非线性关系）；
	\item 对于二维正态分布，独立 $\iff$ 不相关 $\iff$ $\rho = 0$。
\end{enumerate}
\end{myTheorem}

\begin{example}[不相关但不独立的例子]
设 $X \sim N(0,1)$，$Y = X^2$。则
\[\Cov(X, Y) = E(X \cdot X^2) - E(X)E(X^2) = E(X^3) - 0 = 0,\]
故 $X$ 与 $Y$ 不相关。但 $Y$ 完全由 $X$ 决定，显然不独立。
\end{example}
